\chapter{1990年广东卷（文）}

\section{单选题}

\begin{problem}
如果 \(0<\theta<\pi\)，且 \(\cos\theta=-\frac{\sqrt{3}}{2}\)，那么 \(\theta\) 等于（\quad）

\choices
    {\(\frac{5\pi}{6}\)}
    {\(\frac{2\pi}{3}\)}
    {\(\frac{\pi}{3}\)}
    {\(\frac{\pi}{6}\)}
\end{problem}

\begin{answer}
A
\end{answer}

\begin{solution}
在区间 \(0<\theta<\pi\) 内，\(\cos\theta=-\frac{\sqrt{3}}{2}\) 只能在第二象限取得. 由于 \(\cos\frac{5\pi}{6}=-\frac{\sqrt{3}}{2}\)，故 \(\theta=\frac{5\pi}{6}\)，选择 A.
\end{solution}

\begin{problem}
函数 \(f(x)=\sqrt{1+2x}+\sqrt{1-2x}\) 的定义域是（\quad）

\choices
    {\((-\infty,-\frac{1}{2}]\)}
    {\(\left[\frac{1}{2},+\infty\right)\)}
    {\(\left[-\frac{1}{2},+\infty\right)\)}
    {\(\left[-\frac{1}{2},\frac{1}{2}\right]\)}
\end{problem}

\begin{answer}
D
\end{answer}

\begin{solution}
根式有意义要求同时满足 \(1+2x\geqslant0\) 和 \(1-2x\geqslant0\)，即 \(x\geqslant-\frac{1}{2}\) 且 \(x\leqslant\frac{1}{2}\). 取两条件的交集，函数定义域为 \(\left[-\frac{1}{2},\frac{1}{2}\right]\)，选择 D.
\end{solution}

\begin{problem}
椭圆 \(\frac{(x+1)^2}{16}+\frac{y^2}{25}=1\) 的离心率是（\quad）

\choices
    {\(\frac{3}{5}\)}
    {\(\frac{3}{4}\)}
    {\(\frac{\sqrt{41}}{4}\)}
    {\(\frac{\sqrt{41}}{5}\)}
\end{problem}

\begin{answer}
A
\end{answer}

\begin{solution}
将方程写成 \(\frac{y^2}{25}+\frac{(x+1)^2}{16}=1\)，可知半长轴 \(a=5\)，半短轴 \(b=4\). 因而焦半距 \(c=\sqrt{a^2-b^2}=\sqrt{25-16}=3\)，离心率为 \(e=\frac{c}{a}=\frac{3}{5}\)，选择 A.
\end{solution}

\begin{problem}
双曲线 \(\frac{y^2}{16}-\frac{x^2}{9}=1\) 的准线方程是（\quad）

\choices
    {\(x=\pm \frac{16}{\sqrt{7}}\)}
    {\(y=\pm \frac{16}{\sqrt{7}}\)}
    {\(x=\pm \frac{16}{5}\)}
    {\(y=\pm \frac{16}{5}\)}
\end{problem}

\begin{answer}
D
\end{answer}

\begin{solution}
双曲线的标准式为 \(\frac{y^2}{a^2}-\frac{x^2}{b^2}=1\)，所以 \(a=4\)，\(b=3\)，焦半距 \(c=\sqrt{a^2+b^2}=5\)，离心率 \(e=\frac{c}{a}=\frac{5}{4}\). 其准线为 \(y=\pm\frac{a}{e}=\pm\frac{16}{5}\)，选择 D.
\end{solution}

\begin{problem}
\(\cos^2 75^\circ+\cos^2 15^\circ+\cos 75^\circ\cos 15^\circ\) 的值等于（\quad）

\choices
    {\(\frac{\sqrt{6}}{2}\)}
    {\(\frac{5}{4}\)}
    {\(\frac{3}{2}\)}
    {\(1+\frac{\sqrt{3}}{4}\)}
\end{problem}

\begin{answer}
B
\end{answer}

\begin{solution}
因为 \(75^{\circ}+15^{\circ}=90^{\circ}\)，所以 \(\cos^2 75^{\circ}+\cos^2 15^{\circ}=1\). 又
\(\cos75^{\circ}\cos15^{\circ}=\frac{(\sqrt{6}-\sqrt{2})(\sqrt{6}+\sqrt{2})}{16}=\frac{1}{4}\). 原式等于 \(1+\frac14=\frac54\)，选择 B.
\end{solution}

\begin{problem}
如果 \(\tan\theta=2\)，那么 \(\sin^2\theta+\sin\theta\cos\theta+\cos^2\theta=\)（\quad）

\choices
    {\(\frac{7}{3}\)}
    {\(\frac{7}{5}\)}
    {\(\frac{5}{4}\)}
    {\(\frac{5}{3}\)}
\end{problem}

\begin{answer}
B
\end{answer}

\begin{solution}
由 \(\tan\theta=2\)，得 \(\sin\theta=2\cos\theta\)，并且 \(\sin^2\theta+\cos^2\theta=1\). 因此 \(5\cos^2\theta=1\)，从而 \(\sin\theta\cos\theta=2\cos^2\theta=\frac25\). 所求式为 \(1+\sin\theta\cos\theta=1+\frac25=\frac75\)，选择 B.
\end{solution}

\begin{problem}
如果 \(ax+2y+1=0\) 与直线 \(x+y-2=0\) 互相垂直，那么 \(a\) 的值等于（\quad）

\choices
    {\(-2\)}
    {\(-\frac{2}{3}\)}
    {\(-\frac{1}{3}\)}
    {\(1\)}
\end{problem}

\begin{answer}
A
\end{answer}

\begin{solution}
两直线的斜率分别为 \(-\frac{a}{2}\) 和 \(-1\). 垂直时斜率乘积为 \(-1\)，所以 \(\left(-\frac{a}{2}\right)(-1)=-1\)，解得 \(a=-2\)，选择 A.
\end{solution}

\begin{problem}
如果直线 \(l\) 是平面 \(\alpha\) 的斜线，那么平面 \(\alpha\) 内（\quad）

\choices
    {不存在与 \(l\) 平行的直线}
    {不存在与 \(l\) 垂直的直线}
    {与 \(l\) 垂直的直线只有一条}
    {与 \(l\) 平行的直线有无穷多条}
\end{problem}

\begin{answer}
A
\end{answer}

\begin{solution}
斜线 \(l\) 与平面 \(\alpha\) 相交但不在平面内. 若平面 \(\alpha\) 内存在直线与 \(l\) 平行，则由一条直线平行于平面内的直线可推出该直线平行于平面 \(\alpha\)，这与 \(l\) 是斜线相矛盾. 因此不存在平面 \(\alpha\) 内与 \(l\) 平行的直线，选择 A.
\end{solution}

\begin{problem}
某乒乓球队有 \(9\) 名队员，其中 \(2\) 名是种子选手，现要挑选 \(5\) 名队员参加比赛，种子选手都必须在内，那么不同的选法共有（\quad）

\choices
    {\(36\) 种}
    {\(35\) 种}
    {\(84\) 种}
    {\(126\) 种}
\end{problem}

\begin{answer}
B
\end{answer}

\begin{solution}
两名种子选手已经确定入选，还需从另外 \(9-2=7\) 名队员中选出 \(5-2=3\) 名. 不同选法数为 \(\mathrm C_7^3=\frac{7\times6\times5}{3\times2\times1}=35\)，选择 B.
\end{solution}

\begin{problem}
把 \(6\) 个不同的元素排成一排，那么不同的排法共有（\quad）

\choices
    {\(36\) 种}
    {\(120\) 种}
    {\(720\) 种}
    {\(1440\) 种}
\end{problem}

\begin{answer}
C
\end{answer}

\begin{solution}
六个不同元素全部排成一排，第一位有 \(6\) 种选择，第二位有 \(5\) 种选择，依次类推，故排法数为 \(6\times5\times4\times3\times2\times1=6!=720\)，选择 C.
\end{solution}

\begin{problem}
一个圆台的母线长是上、下底面半径的等差中项，且侧面积为 \(8\pi \mathrm{~cm}^2\)，那么母线长为（\quad）

\choices
    {\(4\mathrm{~cm}\)}
    {\(2\sqrt{2}\mathrm{~cm}\)}
    {\(2\mathrm{~cm}\)}
    {\(\sqrt{2}\mathrm{~cm}\)}
\end{problem}

\begin{answer}
C
\end{answer}

\begin{solution}
设圆台上、下底面半径分别为 \(r\) 和 \(R\)，母线长为 \(l\). 由题意 \(l\) 是 \(R\)，\(r\) 的等差中项，所以 \(R+r=2l\). 圆台侧面积为 \(\pi(R+r)l=8\pi\)，于是 \(2\pi l^2=8\pi\)，即 \(l^2=4\). 因为长度为正，得 \(l=2\mathrm{~cm}\)，选择 C.
\end{solution}

\begin{problem}
把复数 \(1+\i\) 对应的向量按顺时针方向旋转 \(\frac{2\pi}{3}\)，所得到的向量对应的复数是（\quad）

\choices
    {\(\frac{1-\sqrt{3}}{2}+\frac{-1+\sqrt{3}}{2}\i\)}
    {\(\frac{1-\sqrt{3}}{2}+\frac{-1-\sqrt{3}}{2}\i\)}
    {\(\frac{-1+\sqrt{3}}{2}+\frac{-1-\sqrt{3}}{2}\i\)}
    {\(\frac{-1+\sqrt{3}}{2}+\frac{1-\sqrt{3}}{2}\i\)}
\end{problem}

\begin{answer}
C
\end{answer}

\begin{solution}
复数 \(1+\i\) 对应向量的模为 \(\sqrt2\)，辐角为 \(\frac{\pi}{4}\). 顺时针旋转 \(\frac{2\pi}{3}\) 后辐角变为 \(\frac{\pi}{4}-\frac{2\pi}{3}=-\frac{5\pi}{12}\). 因而所得复数为
\(\sqrt2\left(\cos\left(-\frac{5\pi}{12}\right)+\i\sin\left(-\frac{5\pi}{12}\right)\right)=\frac{-1+\sqrt3}{2}+\frac{-1-\sqrt3}{2}\i\)，选择 C.
\end{solution}

\begin{problem}
一动点在圆 \(x^2+y^2=1\) 上移动时，它与定点 \((3,0)\) 连线的中点的轨迹方程是（\quad）

\choices
    {\((x+3)^2+y^2=4\)}
    {\((x-3)^2+y^2=1\)}
    {\(\left( x+\frac{3}{2} \right)^2+y^2=\frac{1}{2}\)}
    {\((2x-3)^2+4y^2=1\)}
\end{problem}

\begin{answer}
D
\end{answer}

\begin{solution}
设圆上动点为 \((u,v)\)，线段与定点 \((3,0)\) 的中点为 \((x,y)\)，则 \(u=2x-3\)，\(v=2y\). 代入 \(u^2+v^2=1\)，得 \((2x-3)^2+4y^2=1\)，所以轨迹方程为该式，选择 D.
\end{solution}

\begin{problem}
函数 \(y=\frac{\sin x}{|\sin x|}+\frac{|\cos x|}{\cos x}+\frac{\tan x}{|\tan x|}+\frac{|\cot x|}{\cot x}\) 的值域是（\quad）

\choices
    {\(\{-2,4\}\)}
    {\(\{-2,0,4\}\)}
    {\(\{-2,0,2,4\}\)}
    {\(\{-4,-2,0,4\}\)}
\end{problem}

\begin{answer}
B
\end{answer}

\begin{solution}
函数有意义要求 \(\sin x\ne0\) 且 \(\cos x\ne0\). 令 \(s=\operatorname{sgn}(\sin x)\)，\(c=\operatorname{sgn}(\cos x)\)，则 \(s\)，\(c\in\{-1,1\}\)，且原函数值为 \(s+c+2sc\). 当 \((s,c)=(1,1)\) 时值为 \(4\)，当 \((s,c)=(1,-1)\) 或 \((-1,1)\) 时值为 \(-2\)，当 \((s,c)=(-1,-1)\) 时值为 \(0\). 这些符号组合均可由相应象限中的角取得，所以值域为 \(\{-2,0,4\}\)，选择 B.
\end{solution}

\begin{problem}
复数 \((\sqrt{3}-\i)^3\) 的三角形式是（\quad）

\choices
    {\(8(\cos 90^\circ+\i\sin 90^\circ)\)}
    {\(8(\cos 90^\circ-\i\sin 90^\circ)\)}
    {\(8(\cos 180^\circ+\i\sin 180^\circ)\)}
    {\(8(\cos 270^\circ+\i\sin 270^\circ)\)}
\end{problem}

\begin{answer}
D
\end{answer}

\begin{solution}
\(\sqrt3-\i\) 的模为 \(2\)，辐角为 \(-\frac{\pi}{6}\). 根据棣莫弗定理，\((\sqrt3-\i)^3=8\left(\cos\left(-\frac{\pi}{2}\right)+\i\sin\left(-\frac{\pi}{2}\right)\right)=8(\cos270^{\circ}+\i\sin270^{\circ})\)，选择 D.
\end{solution}

\begin{problem}
如果函数 \(y=(a+1)x-(a-2)x^2\) 是奇函数，那么 \(a\) 的值等于（\quad）

\choices
    {\(-2\)}
    {\(-1\)}
    {\(2\)}
    {\(1\)}
\end{problem}

\begin{answer}
C
\end{answer}

\begin{solution}
要使 \(y=(a+1)x-(a-2)x^2\) 为奇函数，必须有 \(f(-x)=-f(x)\). 其中二次项是偶函数，故其系数必须为零，即 \(-(a-2)=0\)，解得 \(a=2\). 此时 \(y=3x\) 确为奇函数，选择 C.
\end{solution}

\begin{problem}
函数 \(y=(x+4)^2\) 在某区间上是减函数，这区间可以是（\quad）

\choices
    {\(\left[+4,+\infty\right)\)}
    {\(\left(-\infty,+4\right]\)}
    {\(\left[-4,+\infty\right)\)}
    {\(\left(-\infty,-4\right]\)}
\end{problem}

\begin{answer}
D
\end{answer}

\begin{solution}
抛物线 \(y=(x+4)^2\) 的对称轴为 \(x=-4\). 当 \(x\leqslant-4\) 时，\(x\) 增大使 \(|x+4|\) 减小，因此函数值减小，函数在 \(\left(-\infty,-4\right]\) 上是减函数. 选项 D 给出该区间，选择 D.
\end{solution}

\begin{problem}
已知下图（图 4-6-6）是 \(y=2\sin(\omega x+\varphi)\)（\(\left|\varphi\right|<\frac{\pi}{2}\)）的图象，那么（\quad）

\bitmapfigure[max width=0.24\linewidth]{img/1990/guangdong_liberal/q18_fig1.png}

\choices
    {\(\omega=2,\ \varphi=-\frac{\pi}{6}\)}
    {\(\omega=2\)，\(\ \varphi=\frac{\pi}{6}\)}
    {\(\omega=\frac{10}{11}\)，\(\ \varphi=\frac{\pi}{6}\)}
    {\(\omega=\frac{10}{11}\)，\(\ \varphi=-\frac{\pi}{6}\)}
\end{problem}

\begin{answer}
B
\end{answer}

\begin{solution}
由图象在 \(x=0\) 处的函数值为 \(1\)，且曲线向上经过该点，得 \(2\sin\varphi=1\). 结合 \(\left|\varphi\right|<\frac{\pi}{2}\)，可得 \(\varphi=\frac{\pi}{6}\). 图中在 \(x=\frac{11\pi}{12}\) 处向上穿过 \(x\) 轴，因此此处相位为 \(2\pi\)，从而 \(\frac{11\pi}{12}\omega+\frac{\pi}{6}=2\pi\)，解得 \(\omega=2\). 所以选择 B.
\end{solution}

\begin{problem}
如果 \(x\)，\(y\) 是实数，那么“\(xy=0\)”是“\(\left|x+y\right|=\left|x\right|+\left|y\right|\)”的（\quad）

\choices
    {必要而且充分的条件}
    {必要但不充分的条件}
    {充分但不必要的条件}
    {既不充分也不必要的条件}
\end{problem}

\begin{answer}
C
\end{answer}

\begin{solution}
若 \(xy=0\)，则 \(x=0\) 或 \(y=0\)，于是 \(\left|x+y\right|=\left|x\right|+\left|y\right|\)，所以 \(xy=0\) 是后一个条件的充分条件. 反之不成立，例如 \(x=y=1\) 时等式成立但 \(xy=1\ne0\). 因此它是充分但不必要的条件，选择 C.
\end{solution}

\begin{problem}
如果直线 \(y=ax+2\) 与直线 \(y=3x-b\) 关于直线 \(y=x\) 对称，那么（\quad）

\choices
    {\(a=\frac{1}{3},\ b=-6\)}
    {\(a=\frac{1}{3}\)，\(\ b=6\)}
    {\(a=3\)，\(\ b=-2\)}
    {\(a=3\)，\(\ b=6\)}
\end{problem}

\begin{answer}
B
\end{answer}

\begin{solution}
关于直线 \(y=x\) 对称会交换点的横、纵坐标. 直线 \(y=ax+2\) 变为 \(x=ay+2\)，即 \(y=\frac{1}{a}x-\frac{2}{a}\). 与 \(y=3x-b\) 比较斜率和截距，得 \(\frac1a=3\)，\(-\frac2a=-b\)，所以 \(a=\frac13\)，\(b=6\)，选择 B.
\end{solution}

\begin{problem}
如果抛物线 \(y^2=ax\) 的准线方程为 \(x=-1\)，那么它的焦点坐标是（\quad）

\choices
    {\((1,0)\)}
    {\((2,0)\)}
    {\((3,0)\)}
    {\((-1,0)\)}
\end{problem}

\begin{answer}
A
\end{answer}

\begin{solution}
将 \(y^2=ax\) 与标准式 \(y^2=4px\) 比较，得 \(a=4p\). 其准线为 \(x=-p=-1\)，所以 \(p=1\)，焦点为 \((p,0)=(1,0)\)，选择 A.
\end{solution}

\begin{problem}
已知等比数列的公比是 \(2\)，且前 \(4\) 项之和等于 \(1\)，那么前 \(8\) 项之和等于（\quad）

\choices
    {\(21\)}
    {\(19\)}
    {\(17\)}
    {\(15\)}
\end{problem}

\begin{answer}
C
\end{answer}

\begin{solution}
设前四项之和为 \(S_4=1\). 公比为 \(2\)，因此第 \(5\) 至第 \(8\) 项分别是第 \(1\) 至第 \(4\) 项的 \(2^4=16\) 倍，故 \(S_8=S_4+16S_4=17\)，选择 C.
\end{solution}

\begin{problem}
函数 \(y=\tan\frac{x}{2}-\frac{1}{\sin x}\) 的最小正周期是（\quad）

\choices
    {\(\frac{\pi}{2}\)}
    {\(\pi\)}
    {\(\frac{3\pi}{2}\)}
    {\(2\pi\)}
\end{problem}

\begin{answer}
B
\end{answer}

\begin{solution}
函数定义域要求 \(\sin x\ne0\)，这也保证 \(\tan\frac{x}{2}\) 在定义域内有意义. 利用半角公式，令 \(u=\frac{x}{2}\)，则
\[
\tan u-\frac{1}{\sin x}=\frac{\sin u}{\cos u}-\frac{1}{2\sin u\cos u}=\frac{2\sin^2u-1}{2\sin u\cos u}=-\frac{\cos x}{\sin x}=-\cot x
\]
因此原函数与 \(-\cot x\) 在同一定义域上相同，而 \(\cot x\) 的最小正周期为 \(\pi\)，所以选择 B.
\end{solution}

\begin{problem}
以一个三棱柱的顶点为顶点的四面体共有（\quad）

\choices
    {\(30\) 个}
    {\(18\) 个}
    {\(12\) 个}
    {\(6\) 个}
\end{problem}

\begin{answer}
C
\end{answer}

\begin{solution}
三棱柱共有 \(6\) 个顶点，任取四个共有 \(\mathrm C_6^4=15\) 种. 按两个底面分组，若四点中一个底面取 \(3\) 个、另一个底面取 \(1\) 个，共有 \(2\mathrm C_3^3\mathrm C_3^1=6\) 种；若两个底面各取 \(2\) 个，共有 \(\mathrm C_3^2\mathrm C_3^2=9\) 种. 在后一类中，只有上下底面取到同一条边的两个端点时，四点才共面，恰为三条侧面的 \(3\) 组顶点. 因此能组成四面体的四点数为 \(6+9-3=12\)，选择 C.
\end{solution}

\begin{problem}
如果实数 \(x\)，\(y\) 满足等式 \((x-2)^2+y^2=3\)，那么 \(\frac{y}{x}\) 的最大值是（\quad）

\choices
    {\(\frac{1}{2}\)}
    {\(\frac{\sqrt{3}}{3}\)}
    {\(\frac{\sqrt{3}}{2}\)}
    {\(\sqrt{3}\)}
\end{problem}

\begin{answer}
D
\end{answer}

\begin{solution}
圆 \((x-2)^2+y^2=3\) 的圆心为 \((2,0)\)，半径为 \(\sqrt3\). 因为 \(2-\sqrt3>0\)，圆上各点的横坐标均为正，\(\frac{y}{x}\) 是从原点向圆作直线时的斜率. 设切线为 \(y=mx\)，圆心到切线的距离等于半径，故
\[
\frac{2|m|}{\sqrt{m^2+1}}=\sqrt3
\]
平方并化简得 \(4m^2=3(m^2+1)\)，即 \(m^2=3\). 最大斜率为 \(m=\sqrt3\)，选择 D.
\end{solution}

\section{填空题}

\begin{problem}
已知 \(\sin\alpha=\frac{3}{5}\)，\(\alpha\in\left( \frac{\pi}{2},\pi \right)\)，那么 \(\sin \frac{\alpha}{2}\) 的值等于\fillinblank{}.
\end{problem}

\begin{answer}
\(\frac{3\sqrt{10}}{10}\)
\end{answer}

\begin{solution}
因为 \(\alpha\in\left(\frac{\pi}{2},\pi\right)\)，所以 \(\cos\alpha<0\). 由 \(\sin^2\alpha+\cos^2\alpha=1\)，得 \(\cos\alpha=-\sqrt{1-\frac{9}{25}}=-\frac45\). 又 \(0<\frac\alpha2<\frac\pi2\)，故半角公式中的正号适用，\(\sin\frac\alpha2=\sqrt{\frac{1-\cos\alpha}{2}}=\sqrt{\frac{1+\frac45}{2}}=\sqrt{\frac9{10}}=\frac{3\sqrt{10}}{10}\).
\end{solution}

\begin{problem}
\((x-1)-(x-1)^2+(x-1)^3-(x-1)^4\) 的展开式中，\(x^2\) 的系数等于\fillinblank{}.
\end{problem}

\begin{answer}
\(-10\)
\end{answer}

\begin{solution}
逐项考察 \(x^2\) 的系数：\(x-1\) 中没有 \(x^2\) 项，\(-(x-1)^2\) 贡献 \(-1\)，\((x-1)^3\) 贡献 \(-3\)，而 \(-(x-1)^4\) 贡献 \(-6\). 因此总系数为 \(-1-3-6=-10\).
\end{solution}

\begin{problem}
\(\displaystyle\lim_{n\to\infty}\frac{1+\frac{1}{2}+\frac{1}{4}+\cdots+\frac{1}{2^n}}{1+\frac{1}{3}+\frac{1}{9}+\cdots+\frac{1}{3^n}}=\)\fillinblank{}.
\end{problem}

\begin{answer}
\(\frac{4}{3}\)
\end{answer}

\begin{solution}
分子和分母分别是等比数列的前 \(n+1\) 项和，所以
\[
1+\frac12+\cdots+\frac1{2^n}=\frac{1-2^{-(n+1)}}{1-\frac12}\longrightarrow2
\]
又
\[
1+\frac13+\cdots+\frac1{3^n}=\frac{1-3^{-(n+1)}}{1-\frac13}\longrightarrow\frac32
\]
分母的极限非零，故所求极限为 \(\frac{2}{3/2}=\frac43\).
\end{solution}

\begin{problem}
如果三棱锥每条侧棱长都是 \(a\)，且每两条侧棱所成的角都是 \(60^\circ\)，那么这个三棱锥的高的长度是\fillinblank{}.
\end{problem}

\begin{answer}
\(\frac{\sqrt{6}}{3}a\)
\end{answer}

\begin{solution}
设三条侧棱对应的向量为 \(\overrightarrow{SA}=\bs{u}\)，\(\overrightarrow{SB}=\bs{v}\)，\(\overrightarrow{SC}=\bs{w}\). 由题意 \(|\bs{u}|=|\bs{v}|=|\bs{w}|=a\)，且任意两向量的夹角为 \(60^\circ\)，所以任意两者的数量积均为 \(\frac{a^2}{2}\). 设 \(G\) 为底面三角形 \(ABC\) 的重心，则 \(\overrightarrow{SG}=\frac{\bs{u}+\bs{v}+\bs{w}}{3}\). 对底面内的两条方向向量 \(\bs{v}-\bs{u}\)、\(\bs{w}-\bs{u}\)，有 \((\bs{u}+\bs{v}+\bs{w})\cdot(\bs{v}-\bs{u})=0\) 且 \((\bs{u}+\bs{v}+\bs{w})\cdot(\bs{w}-\bs{u})=0\)，故 \(SG\perp\) 面 \(ABC\)，三棱锥的高为 \(|SG|\). 又
\[
|\bs{u}+\bs{v}+\bs{w}|^2=3a^2+6\cdot\frac{a^2}{2}=6a^2
\]
所以 \(SG=\frac{\sqrt6}{3}a\).
\end{solution}

\begin{problem}
如果 \(z_1\)，\(z_2\) 是复数，且 \(\left| z_1 \right|=3\)，\(\left| z_2 \right|=4\)，\(\left| z_1-z_2 \right|=5\)，那么 \(\left| z_1+z_2 \right|\) 的值是\fillinblank{}.
\end{problem}

\begin{answer}
\(5\)
\end{answer}

\begin{solution}
利用模的平方公式，\(\left|z_1-z_2\right|^2=|z_1|^2+|z_2|^2-2\operatorname{Re}(z_1\overline{z_2})\). 代入已知值得 \(25=9+16-2\operatorname{Re}(z_1\overline{z_2})\)，所以 \(\operatorname{Re}(z_1\overline{z_2})=0\). 因而 \(\left|z_1+z_2\right|^2=|z_1|^2+|z_2|^2+2\operatorname{Re}(z_1\overline{z_2})=9+16=25\). 模是非负数，故 \(\left|z_1+z_2\right|=5\).
\end{solution}

\section{解答题}

\begin{problem}
已知 \(f(x)=\log_3\left( 4+3x-x^2 \right)-\log_3\left( 2x-1 \right)\).

\begin{enumerate}
\item 求函数 \(f(x)\) 的定义域；
\item 解不等式 \(f(x)>\log_3 2\).
\end{enumerate}
\end{problem}

\begin{solution}
\begin{enumerate}
\item 对数的真数必须为正，故定义域由不等式组
\[
\begin{cases}
4+3x-x^2>0,\\
2x-1>0.
\end{cases}
\]
第一式可化为 \((x+1)(x-4)<0\)，所以 \(-1<x<4\)；第二式给出 \(x>\frac12\). 取交集得 \(\frac12<x<4\)，所以函数 \(f(x)\) 的定义域是 \(\left\{x\mid\frac12<x<4\right\}\).

\item 在上述定义域内，由于对数的底数 \(3>1\)，不等式等价于
\[
\log_3\frac{4+3x-x^2}{2x-1}>\log_3 2
\]
即
\[
\begin{cases}
\frac{1}{2}<x<4,\\
4+3x-x^2>2(2x-1).
\end{cases}
\]
第二个不等式化为 \(x^2+x-6<0\)，即 \(-3<x<2\). 与定义域取交集，得 \(\frac12<x<2\). 因此不等式的解集是 \(\left\{x\mid\frac12<x<2\right\}\).
\end{enumerate}
\end{solution}

\begin{problem}
如图，在三棱锥 \(S-ABC\) 中，\(SA\perp\) 底面 \(ABC\)，\(AB\perp BC\)，\(DE\) 垂直平分 \(SC\)，且分别交 \(AC\)，\(SC\) 于 \(D\)，\(E\)，又 \(SA=AB=a\)，\(BC=\sqrt{2}a\).

\begin{enumerate}
\item 求证：\(SC\perp\) 面 \(BDE\)；
\item 求面 \(BDE\) 与面 \(BDC\) 所成的二面角的大小.
\end{enumerate}

\FigureLayoutDeclare{img/1990/guangdong_liberal/q2_fig1.png}{6.8cm}{65bp 0bp 50bp 1bp}
\bitmapfigure[max width=0.24\linewidth]{img/1990/guangdong_liberal/q2_fig1.png}
\end{problem}

\begin{solution}
\begin{enumerate}
\item 建立如图所示的直角坐标系，使原点为 \(A\)，三个坐标轴分别沿 \(AB\)，\(BC\)，\(SA\) 方向. 则 \(A(0,0,0)\)，\(B(a,0,0)\)，\(C(a,\sqrt{2}a,0)\)，\(S(0,0,a)\). 设 \(D=(ta,\sqrt{2}ta,0)\). 由 \(DE\) 垂直平分 \(SC\)，得 \(DS=DC\)，于是 \(3t^2a^2+a^2=3(1-t)^2a^2\)，解得 \(t=\frac13\). 因此 \(D\left(\frac a3,\frac{\sqrt{2}a}{3},0\right)\)，\(E\left(\frac a2,\frac{\sqrt{2}a}{2},\frac a2\right)\). 又 \(\overrightarrow{SC}=(a,\sqrt{2}a,-a)\)，\(\overrightarrow{DB}=\left(\frac{2a}{3},-\frac{\sqrt{2}a}{3},0\right)\)，\(\overrightarrow{DE}=\left(\frac a6,\frac{\sqrt{2}a}{6},\frac a2\right)\). 所以 \(\overrightarrow{SC}\cdot\overrightarrow{DB}=0\)，\(\overrightarrow{SC}\cdot\overrightarrow{DE}=0\).
由于 \(DB\) 与 \(DE\) 是平面 \(BDE\) 内的相交直线，故 \(SC\perp\) 面 \(BDE\).

\item 取平面 \(BDE\) 的一个法向量 \(\bs{n}_1=\overrightarrow{DB}\times\overrightarrow{DE}=\frac{a^2}{6}(-\sqrt{2},-2,\sqrt{2})\)，平面 \(BDC\) 的一个法向量为 \(\bs{n}_2=(0,0,1)\). 设二面角为 \(\theta\)，则 \(\cos\theta=\frac{|\bs{n}_1\cdot\bs{n}_2|}{|\bs{n}_1||\bs{n}_2|}=\frac12\).
因此二面角的大小为 \(\theta=\frac{\pi}{3}\).
\end{enumerate}
\end{solution}

\begin{problem}
已知椭圆的焦点为 \(F_1(-1,0)\) 和 \(F_2(1,0)\)，直线 \(x=4\) 是椭圆的一条准线.

\begin{enumerate}
\item 求椭圆的方程；
\item 又设点 \(P\) 在椭圆上，且 \(\left| PF_1 \right|-\left| PF_2 \right|=1\)，求 \(\cos\angle F_1PF_2\) 的值.
\end{enumerate}
\end{problem}

\begin{solution}
\begin{enumerate}
\item 设椭圆方程为 \(\frac{x^2}{a^2}+\frac{y^2}{b^2}=1\)，焦半距为 \(c\). 由焦点坐标知 \(c=1\). 椭圆的右准线方程为 \(x=\frac{a^2}{c}\)，所以由 \(x=4\) 得 \(a^2=4\)，即 \(a=2\)，进而 \(b^2=a^2-c^2=3\). 所以椭圆方程为 \(\frac{x^2}{4}+\frac{y^2}{3}=1\).

\item 设 \(d_1=|PF_1|\)，\(d_2=|PF_2|\). 由椭圆定义及题设，得 \(d_1+d_2=2a=4\)，\(d_1-d_2=1\)，从而 \(d_1=\frac52\)，\(d_2=\frac32\). 在三角形 \(F_1PF_2\) 中，\(F_1F_2=2\)，由余弦定理
\[
\cos\angle F_1PF_2=\frac{d_1^2+d_2^2-F_1F_2^2}{2d_1d_2}=\frac{\frac{25}{4}+\frac94-4}{2\times\frac52\times\frac32}=\frac35
\]
所以 \(\cos\angle F_1PF_2=\frac35\).
\end{enumerate}
\end{solution}

\begin{problem}
已知数列 \(\{a_n\}\) 的前 \(n\) 项和的公式 \(S_n=\frac{\pi}{12}(2n^2+n)\).

\begin{enumerate}
\item 求证 \(\{a_n\}\) 是等差数列，并求出它的首项的公差；
\item 记 \(b_n=\sin a_n\cdot \sin a_{n+1}\cdot \sin a_{n+2}\)，求证：对任何自然数 \(n\)，都有 \(b_n=\frac{\sqrt{2}}{8}(-1)^{n-1}\).
\end{enumerate}
\end{problem}

\begin{solution}
\begin{enumerate}
\item 由 \(a_1=S_1=\frac\pi4\)，且当 \(n\geqslant2\) 时
\(a_n=S_n-S_{n-1}=\frac\pi{12}(4n-1)\)，
故对一切自然数 \(n\)，都有 \(a_n=\frac\pi{12}(4n-1)\). 因此
\(a_{n+1}-a_n=\frac\pi3\)，
所以 \(\{a_n\}\) 是等差数列，首项为 \(\frac\pi4\)，公差为 \(\frac\pi3\).

\item 由上式知 \(a_{n+3}=a_n+\pi\)，且 \(4n-1\) 不是 \(12\) 的倍数，所以 \(\sin a_n\neq0\). 先算得
\(b_1=\sin\frac\pi4\sin\frac{7\pi}{12}\sin\frac{11\pi}{12}=\frac{\sqrt{2}}8\).
又
\(\frac{b_{n+1}}{b_n}=\frac{\sin a_{n+3}}{\sin a_n}=-1\)，
所以 \(\{b_n\}\) 是首项为 \(\frac{\sqrt{2}}8\)，公比为 \(-1\) 的等比数列，从而对任何自然数 \(n\)，都有
\(b_n=\frac{\sqrt{2}}8(-1)^{n-1}\).
\end{enumerate}
\end{solution}
